%% LyX 2.4.4 created this file.  For more info, see https://www.lyx.org/.
%% Do not edit unless you really know what you are doing.
\documentclass[oneside,english]{amsart}
\usepackage[T1]{fontenc}
\usepackage[latin9]{inputenc}
\synctex=-1
\usepackage{amsthm}

\makeatletter
%%%%%%%%%%%%%%%%%%%%%%%%%%%%%% Textclass specific LaTeX commands.
\newlength{\lyxlabelwidth}      % auxiliary length 
\theoremstyle{plain}
\newtheorem*{thm*}{\protect\theoremname}
\theoremstyle{remark}
\newtheorem*{claim*}{\protect\claimname}

%%%%%%%%%%%%%%%%%%%%%%%%%%%%%% User specified LaTeX commands.
\date{}
\usepackage{commath}

\makeatother

\usepackage{babel}
\providecommand{\claimname}{Claim}
\providecommand{\theoremname}{Theorem}

\begin{document}
\begin{thm*}
Here is some theorem.
\end{thm*}
\begin{proof}
This is the beginning of the proof.
\begin{claim*}
Here is what a claim generally looks like when stated inside a proof.
There is no indentation on the first line or on subsequent lines when
the statement of the claim goes on past a line break.
\end{claim*}
This is the next paragraph of the proof, and this is where we both
prove the claim and complete the proof.
\end{proof}
Here is ordinary text, following the proof.
\end{document}
